\chapter{Introducción}

\noindent Los \textbf{\glosref{mercado-digital}{mercados digitales}} son entornos en línea en los que las personas usuarias intercambian bienes o servicios mediante \textbf{\glosref{plataforma}{plataformas}}. Una plataforma administra las reglas del intercambio y la información necesaria para participar en él. Cuando los bienes están asociados a una cuenta o a un servicio digital y pueden utilizarse, coleccionarse o intercambiarse, se denominan \textbf{\glosref{activo-digital}{activos digitales}}.

La gestión de activos digitales plantea decisiones sucesivas de compra, venta, mantenimiento o descarte. Estas decisiones se denominan \textbf{\glosref{operacion}{operaciones}}. Cada operación puede modificar la \textbf{\glosref{cartera}{cartera}}, formada por el saldo disponible y las \textbf{\glosref{posicion-abierta}{posiciones abiertas}}, es decir, los activos digitales adquiridos que aún no se han vendido. Por ello, el objetivo no es solo identificar una diferencia de precios, sino evaluar decisiones que afectan a la evolución de la cartera con el fin de generar un beneficio económico.

El \textbf{\glosref{aprendizaje-refuerzo}{aprendizaje por refuerzo}} (\acroref{RL}{RL}), del inglés \emph{Reinforcement Learning}, es un paradigma para abordar problemas de decisión secuencial. En este marco, un \textbf{\glosref{agente}{agente}} interactúa con un \textbf{\glosref{entorno}{entorno}}, selecciona acciones y recibe \textbf{\glosref{recompensa}{recompensas}}. A diferencia de un modelo que solo estima un resultado futuro, RL permite estudiar cómo una decisión actual condiciona las decisiones posteriores.

El \textbf{\glosref{aprendizaje-refuerzo-multiagente}{aprendizaje por refuerzo multiagente}} (\acroref{MARL}{MARL}), del inglés \emph{Multi-Agent Reinforcement Learning}, amplía este marco cuando varios agentes interactúan con el mismo entorno. Los agentes pueden cooperar, competir o actuar de forma independiente. La coordinación puede ser centralizada, cuando una única entidad toma la decisión conjunta, o descentralizada, cuando cada agente decide desde su propia información y responsabilidad.

\section{Motivación}

Los mercados de activos digitales en videojuegos permiten que las personas usuarias compren, vendan e intercambien objetos asociados a un videojuego. El valor de estos activos digitales puede depender de su utilidad (por ejemplo, aumenta la vida de un personaje), apariencia, \textbf{\glosref{rareza}{rareza}} (la dificultad para obtener un activo digital), \textbf{\glosref{disponibilidad}{disponibilidad}} (el número de unidades en venta) y \textbf{\glosref{demanda}{demanda}} (el número de personas que quieren adquirirlo). La literatura sobre bienes virtuales identifica dimensiones funcionales, estéticas y sociales como determinantes de ese valor \cite{lehdonvirta2009virtual}. Por ello, una misma operación puede presentar un resultado económico diferente según el momento y la plataforma en los que se realice.

La motivación principal de este trabajo es apoyar la identificación de operaciones que puedan generar un \textbf{\glosref{beneficio-neto}{beneficio neto}} (el resultado después de descontar el precio de compra, las comisiones y otros costes). Una decisión acertada puede aportar un resultado económico favorable a la persona usuaria. Sin embargo, ese resultado no está garantizado y comparar dos precios no es suficiente para decidir. Un mismo activo digital puede aparecer en varias plataformas con precios, divisas, comisiones y condiciones de intercambio distintas. La decisión exige valorar cuánto se pagaría al comprar, cuánto se recibiría al vender y si se aplica un \textbf{\glosref{trade-hold}{\textit{trade hold}}} (bloqueo temporal que impide vender o transferir el activo adquirido).

\section{Planteamiento del problema}

Las decisiones de compra y venta se suceden en el tiempo. Cada compra utiliza una parte del saldo disponible de la cartera digital. Mientras el activo digital adquirido permanezca en la cartera y no se pueda vender, ese importe no puede utilizarse para adquirir otro activo digital. Por ello, el problema no consiste solo en comparar precios, sino también en decidir cómo emplear un saldo limitado entre distintas posibilidades de compra.

Las plataformas no siempre emplean la misma moneda, el mismo nombre ni la misma fecha de observación para un activo digital. Antes de comparar precios, es necesario unificar esa información.

Cuando se analizan conjuntamente una compra y una venta posterior, se habla de una \textbf{\glosref{operacion-compraventa}{operación de compraventa}}. El punto de partida es un procedimiento manual en el que, para cada activo digital, se anotan el precio de compra, el precio de venta, la conversión de moneda, las comisiones, la fecha de desbloqueo y el \textbf{\glosref{capital-bloqueado}{capital bloqueado}} (el saldo que no puede utilizarse para nuevas compras). Esta primera implementación manual de donde partimos permite evaluar casos concretos, pero obliga a repetir los cálculos y a comprobar los datos de cada activo digital.

Con la información mencionada anteriormente, el sistema debe decidir qué operaciones conviene realizar, mantener o descartar. Para justificar esa decisión, también debe considerar el \textbf{\glosref{volumen-intercambios}{volumen de intercambios}} (la cantidad de unidades vendidas durante un día, una semana o un mes), las reglas económicas y de intercambio aplicables y el estado de la cartera. La pregunta que debe responder es si la posible operación merece ser estudiada con criterios consistentes.

\section{Objetivos}

El objetivo principal de este \acroref{TFM}{TFM} es diseñar y evaluar, mediante simulación, una arquitectura multiagente que apoye la identificación y evaluación de operaciones orientadas a generar un beneficio neto en mercados de activos digitales. El caso de estudio es la economía virtual del videojuego \textit{Counter-Strike 2}.

Los objetivos específicos se agrupan en cuatro bloques:

\begin{itemize}
  \item \textbf{Análisis del caso de estudio}: analizar el mercado de activos digitales de \textit{Counter-Strike 2}, sus plataformas, restricciones, comisiones, divisas, volumen de intercambios y efectos del \textit{trade hold}.
  \item \textbf{Construcción de datos}: construir un flujo de adquisición y preparación de datos procedentes de diferentes plataformas de mercado y fuentes manuales, manteniendo la trazabilidad de la fuente, la fecha, la moneda y la calidad del dato.
  \item \textbf{Simulación y estimación}: desarrollar un simulador que incorpore las reglas económicas del procedimiento manual, construir conjuntos de datos temporales y obtener estimaciones de beneficio que puedan incorporarse al entorno de decisión.
  \item \textbf{Coordinación multiagente}: formular un entorno de aprendizaje por refuerzo con el estado de la cartera, las acciones y las recompensas definidos. Sobre esa base, distribuir los roles de detección, decisión y control de riesgo entre agentes especializados y evaluar su funcionamiento con métricas de \textbf{\glosref{rentabilidad}{rentabilidad}}, riesgo, \textbf{\glosref{exposicion}{exposición}} y capital bloqueado.
\end{itemize}

\section{Estructura del documento}

La memoria se estructura de forma secuencial siguiendo las fases principales del proyecto. A continuación, se detalla el contenido de cada capítulo:

\begin{itemize}
  \item \textbf{Capítulo 1:} introduce el contexto del proyecto, la motivación, el problema abordado, la solución propuesta, los objetivos, el alcance y la estructura.
  \item \textbf{Capítulo 2:} presenta el estado del arte sobre mercados de activos digitales, plataformas de intercambio, aprendizaje por refuerzo, coordinación multiagente y toma de decisiones en carteras.
  \item \textbf{Capítulo 3:} analiza los conceptos operativos del mercado de activos digitales en videojuegos y justifica el uso de \textit{Counter-Strike 2} como caso de estudio.
  \item \textbf{Capítulo 4:} concreta el modelo de decisión multiagente aplicado al caso de estudio y presenta la arquitectura propuesta. Se describen el estado, las observaciones, las acciones, la transición simulada, los agentes, la función de recompensa y las restricciones de riesgo.
  \item \textbf{Capítulo 5:} expone el entorno de datos y simulación. Se detallan fuentes de adquisición, persistencia, características utilizadas, simulador de cartera, \acroref{OCR}{OCR} e interfaz de consulta.
  \item \textbf{Capítulo 6:} describe la implementación y operación del sistema. Presenta la organización del repositorio, el scraping, la persistencia, la consola local y la configuración reproducible.
  \item \textbf{Capítulo 7:} recoge la experimentación realizada. Incluye migración de reglas económicas, construcción de datasets supervisados, evaluación de modelos tabulares, pruebas técnicas y de rendimiento, dataset de trading, estado del entrenamiento \acroref{MARL}{MARL}, dificultades y pivotaciones detectadas.
  \item \textbf{Capítulo 8:} presenta conclusiones principales, revisa el cumplimiento de objetivos y plantea líneas de trabajo necesarias para completar la validación experimental.
  \item \textbf{Anexos:} reúnen información complementaria sobre estructura del repositorio, configuración del entorno, comandos de ejecución, parámetros de entrenamiento y modelo de datos.
\end{itemize}
